\section*{अध्याय \ref{ch:g11:organic-skeletons} --- कार्बनिक अणु: कंकाल और नाम}

\begin{solution}{exo:g11:organic-skeletons:1}
दोनों में छह कार्बन \omterm{def:g7:atoms-and-molecules:atom}{परमाणु} (सिरे और कोने) हैं और \ce{C6H14}: हेक्सेन और
2-मेथिलपेंटेन \omterm{def:g11:organic-skeletons:isomer}{समावयवी} हैं।
\end{solution}

\begin{solution}{exo:g11:organic-skeletons:2}
प्रोपेन; हेक्सेन; 2-मेथिलब्यूटेन।
\end{solution}

\begin{solution}{exo:g11:organic-skeletons:3}
\ce{CH3-CH2-CH3}; \chemfig{CH_3-CH(-[2]CH_3)-CH_3};
\chemfig{H-C(-[2]H)(-[6]H)-C(-[2]H)(-[6]H)-H}।
\end{solution}

\begin{solution}{exo:g11:organic-skeletons:4}
पहला, 3-मेथिलपेंटेन: 6 कार्बन \omterm{def:g7:atoms-and-molecules:atom}{परमाणु}; उसके \ce{CH3} सिरे (तीन)
9 हाइड्रोजन \omterm{def:g7:atoms-and-molecules:atom}{परमाणु} रखते हैं, उसके दो \ce{CH2} कोने 4 और उसका
\ce{CH} शाखा-बिंदु 1: \ce{C6H14}। साइक्लोहेक्सेन: 6 कोने, हर एक
\ce{CH2}: \ce{C6H12}।
\end{solution}

\begin{solution}{exo:g11:organic-skeletons:5}
\ce{C8H18}। $2n + 2 = 20$ से $n = 9$: नोनेन, \ce{C9H20}।
\end{solution}

\begin{solution}{exo:g11:organic-skeletons:6}
\setchemfig{atom sep=1.6em}
3-मेथिलहेक्सेन \chemfig{-[:30]-[:-30]-[:30](-[:90])-[:-30]-[:30]-[:-30]},
जिसका सूत्र \ce{C7H16} है; 2,2-डाइमेथिलब्यूटेन \chemfig{-[:0](-[:90])(-[:-90])-[:0]-[:30]},
जिसका सूत्र \ce{C6H14} है; 3-एथिलपेंटेन
\chemfig{-[:30]-[:-30](-[:-90]-[:-30])-[:30]-[:-30]}, जिसका सूत्र \ce{C7H16} है।
\end{solution}

\begin{solution}{exo:g11:organic-skeletons:7}
पेंटेन, 2-मेथिलब्यूटेन और 2,2-डाइमेथिलप्रोपेन, जो अध्याय के चित्र में
बने हैं।
\end{solution}

\begin{solution}{exo:g11:organic-skeletons:8}
``2-एथिलब्यूटेन'' में एथिल समूह से होकर जाने वाली सबसे लंबी श्रृंखला पाँच कार्बनों
की है: 3-मेथिलपेंटेन। ``4-मेथिलपेंटेन'' पर ग़लत सिरे से क्रमांक दिए गए हैं:
2-मेथिलपेंटेन। ``1-मेथिलब्यूटेन'' बस पाँच कार्बनों की एक
श्रृंखला है: पेंटेन।
\end{solution}

\begin{solution}{exo:g11:organic-skeletons:9}
2,2-डाइमेथिलप्रोपेन (\qty{9.5}{\celsius}) < 2-मेथिलब्यूटेन
(\qty{27.8}{\celsius}) < पेंटेन (\qty{36.1}{\celsius})। वही \omterm{def:g7:atoms-and-molecules:atom}{परमाणु}
उत्तरोत्तर अधिक सघन ढंग से सजे हैं: \omterm{def:g7:atoms-and-molecules:molecule}{अणुओं} के बीच कम संपर्क,
दुर्बल \omterm{def:g11:polarity-and-cohesion:van-der-waals}{वान्डरवाल्स अन्योन्यक्रियाएँ}, कम क्वथन ताप।
\end{solution}

\begin{solution}{exo:g11:organic-skeletons:10}
श्रृंखला को वलय में बंद करने से एक और \ce{C-C} आबंध बनता है, जो
दो \ce{C-H} आबंधों (हर सिरे पर एक) की जगह लेता है। \omterm{def:g11:organic-skeletons:formulas}{आण्विक सूत्र} अलग:
\omterm{def:g11:organic-skeletons:isomer}{समावयवी} नहीं। साइक्लोहेक्सेन \omterm{def:g11:organic-skeletons:alkane}{ऐल्केन} नहीं है (उसका कंकाल एक
वलय है; उसका सूत्र \ce{C_{n}H_{2n+2}} नहीं है)।
\end{solution}

\begin{solution}{exo:g11:organic-skeletons:11}
ब्यूटेन और 2-मेथिलप्रोपेन (\ce{C4H10})। प्रोपेन \ce{C3H8} है और
साइक्लोब्यूटेन \ce{C4H8}।
\end{solution}

\begin{solution}{exo:g11:organic-skeletons:12}
\setchemfig{atom sep=1.5em}
\begin{center}
\small
\begin{tabular}{ccc}
\chemfig{-[:30]-[:-30]-[:30]-[:-30]-[:30]} &
\chemfig{-[:30](-[:90])-[:-30]-[:30]-[:-30]} &
\chemfig{-[:30]-[:-30](-[:-90])-[:30]-[:-30]} \\
हेक्सेन & 2-मेथिलपेंटेन & 3-मेथिलपेंटेन \\[6pt]
\chemfig{-[:0](-[:90])(-[:-90])-[:0]-[:30]} &
\chemfig{-[:30](-[:90])-[:-30](-[:-90])-[:30]} & \\
2,2-डाइमेथिलब्यूटेन & 2,3-डाइमेथिलब्यूटेन &
\end{tabular}
\end{center}
\end{solution}

\begin{solution}{exo:g11:organic-skeletons:13}
दो लंबी टेढ़ी-मेढ़ी रेखाएँ अपनी पूरी लंबाई पर अगल-बगल लेट सकती हैं, जिनमें
बहुत-से \omterm{def:g7:atoms-and-molecules:atom}{परमाणु} संपर्क में रहते हैं; दो गेंदें केवल एक छोटे-से क्षेत्र पर छूती हैं।
\omterm{def:g11:polarity-and-cohesion:van-der-waals}{वान्डरवाल्स अन्योन्यक्रियाएँ} संपर्क वाले \omterm{def:g7:atoms-and-molecules:atom}{परमाणुओं} के बीच काम करती हैं: वे
रैखिक \omterm{def:g7:atoms-and-molecules:molecule}{अणुओं} के बीच अधिक होती हैं, जिन्हें अलग करने के लिए ऊँचा ताप चाहिए।
\end{solution}

\begin{solution}{exo:g11:organic-skeletons:14}
सबसे लंबी श्रृंखला में छह कार्बन हैं; बाएँ से क्रमांक देने पर
प्रतिस्थापी 2, 4, 4 पर हैं, दाएँ से 3, 3, 5 पर; पहला
अंतर (3 के मुक़ाबले 2) फ़ैसला करता है: 2,4,4-ट्राइमेथिलहेक्सेन।
\end{solution}

\begin{solution}{exo:g11:organic-skeletons:15}
\setchemfig{atom sep=1.6em}
\chemfig{-[:0](-[:90])(-[:-90])-[:-30]-[:30](-[:90])-[:-30]}: पाँच कार्बनों की
एक श्रृंखला, जिसके कार्बन 2 पर दो मेथिल समूह और कार्बन 4 पर एक है,
कुल आठ कार्बन \omterm{def:g7:atoms-and-molecules:atom}{परमाणु}; $2 \times 8 + 2 = 18$ हाइड्रोजन \omterm{def:g7:atoms-and-molecules:atom}{परमाणु},
\ce{C8H18}, ऑक्टेन का सूत्र: वे \omterm{def:g11:organic-skeletons:isomer}{समावयवी} हैं। \omterm{def:g11:organic-skeletons:formulas}{अर्ध-संरचना
सूत्र}: \chemfig{CH_3-C(-[2]CH_3)(-[6]CH_3)-CH_2-CH(-[2]CH_3)-CH_3}।
\end{solution}

\begin{solution}{pb:g11:organic-skeletons:1}
\setchemfig{atom sep=1.5em}
\textbf{1.} \ce{C4H10}: $n = 4$ के साथ, $2n + 2 = 10$।

\textbf{2.} \chemfig{H-C(-[2]H)(-[6]H)-C(-[2]H)(-[6]H)-C(-[2]H)(-[6]H)-C(-[2]H)(-[6]H)-H}।

\textbf{3.} \ce{CH3-CH2-CH2-CH3}; \chemfig{-[:30]-[:-30]-[:30]}।

\textbf{4.} \chemfig{H-C(-[2]H)(-[6]H)-C(-[6]H)(-[2,1.5]C(-[1]H)(-[3]H)(-[2]H))-C(-[2]H)(-[6]H)-H};
\chemfig{CH_3-CH(-[2]CH_3)-CH_3}; \chemfig{-[:30](-[:90])-[:-30]}। उसमें
वही \omterm{def:g7:atoms-and-molecules:atom}{परमाणु} हैं, चार कार्बन और दस हाइड्रोजन, जो
अलग क्रम में आबंधित हैं।

\textbf{5.} हाँ: \omterm{def:g11:organic-skeletons:isomer}{संरचनात्मक समावयवी} (अलग कंकाल)।

\textbf{6.} तीनों: उनके क्वथन ताप \qty{20}{\celsius} से
नीचे हैं।

\textbf{7.} \qty{-5}{\celsius} पर, अपने क्वथन ताप से नीचे, ब्यूटेन
सामान्य दाब पर द्रव बनी रहती है और थोड़ी ही गैस देती है।

\textbf{8.} प्रोपेन (\qty{-42}{\celsius}) और 2-मेथिलप्रोपेन
(\qty{-11.7}{\celsius}) सर्दियों के तापों से नीचे उबलती हैं: वे ठंड में भी
भरपूर गैस देती हैं।

\textbf{9.} 2-मेथिलप्रोपेन शाखित है, अधिक सघन: दुर्बल
\omterm{def:g11:polarity-and-cohesion:van-der-waals}{वान्डरवाल्स अन्योन्यक्रियाएँ}।

\textbf{10.} तीन कार्बन \omterm{def:g7:atoms-and-molecules:atom}{परमाणुओं} के साथ एकमात्र संभव कंकाल श्रृंखला
\ce{C-C-C} है: किसी भी कार्बन को कहीं और ले जाने पर वही श्रृंखला मिलती है।

\textbf{11.} 2-मेथिलब्यूटेन।

\textbf{12.} श्रृंखला पर उस सिरे से क्रमांक देने चाहिए जो प्रतिस्थापी के
अधिक पास है: 2, 3 नहीं।

\textbf{13.} 2, 3 और 5।

\textbf{14.} हेप्टेन, \chemfig{-[:30]-[:-30]-[:30]-[:-30]-[:30]-[:-30]}।

\textbf{15.} ये दो हैं: 2-मेथिलहेक्सेन \chemfig{-[:30](-[:90])-[:-30]-[:30]-[:-30]-[:30]}
और 3-मेथिलहेक्सेन \chemfig{-[:30]-[:-30](-[:-90])-[:30]-[:-30]-[:30]}।
``4-मेथिलहेक्सेन'' ग़लत सिरे से क्रमांकित 3-मेथिलहेक्सेन है।

\textbf{16.} 2,2-डाइमेथिलपेंटेन, 2,3-डाइमेथिलपेंटेन,
2,4-डाइमेथिलपेंटेन और 3,3-डाइमेथिलपेंटेन।

\textbf{17.} हाँ, 3-एथिलपेंटेन। कार्बन 2 पर एथिल समूह उससे होकर एक छह-कार्बन
श्रृंखला बना देता: तब \omterm{def:g7:atoms-and-molecules:molecule}{अणु} 3-मेथिलहेक्सेन होता।

\textbf{18.} 2,2,3-ट्राइमेथिलब्यूटेन,
\chemfig{-[:0](-[:90])(-[:-90])-[:0](-[:90])-[:0]}।

\textbf{19.} $1 + 2 + 4 + 1 + 1 = 9$।

\textbf{20.} हेप्टेन के, स्वयं उसे मिलाकर, 9 \omterm{def:g11:organic-skeletons:isomer}{संरचनात्मक समावयवी} हैं।
\end{solution}
